\documentclass[11pt, a4paper]{article}

\usepackage[utf8]{inputenc}
\usepackage{amsmath, amssymb, amsthm}
\usepackage{hyperref}
\usepackage{geometry}
\usepackage{graphicx}
\usepackage{algorithmic}
\usepackage{xcolor}

\geometry{margin=1in}

\title{\textbf{Architectural Evolution of the "UTGS" Engine: \\ 
A Decoupled, Multi-Backend, Data-Driven \\
Game Development Framework (2001--2006)
}}
\author{Sonia Kolasinska}
\date{March 2026 (Retrospective)}

\begin{document}

\maketitle

\setlength{\parskip}{1em}
\setlength{\parindent}{0pt}

\begin{abstract}
This paper details the technical design of a vertically integrated C++ framework developed between 2001 and 2006. The system pioneered a custom functional language, \textbf{KooLiXP}, utilizing an optimizing compiler to transform static XML into high-performance executable ASTs. A critical innovation of the engine was its \textbf{Layered Workspace Architecture}, which utilized a priority-based Z-order to manage complex UI states (Desktop, Popups, Menus, and Hints). By decoupling application logic from execution via KooLiXP and employing a modular Hardware Abstraction Layer (HAL) for OpenGL, DirectDraw, and GDI, the framework enabled rapid development of high-fidelity applications with robust system-level stability.
\end{abstract}

\section{Introduction}
Early 2000s graphics programming was a landscape of hardware fragmentation. Developers were forced to choose between raw performance and modularity. The "UTGS" ("Useless Temporary Game Solutions") engine resolved this tension by implementing a "Full-Stack" abstraction. The architecture provided an optimizing compiler for the \textbf{KooLiXP} language, a native plugin system (\textbf{Dynamo}), and a sophisticated UI compositor (\textbf{Workspace}) that managed visual depth and input routing across distinct logical layers.

\section{The KooLiXP Programming Language and Virtual Machine}
The defining cornerstone of the engine is its proprietary scripting language, \textbf{KooLiXP} (derived from "Koolas", the author's surname, and "LiXP", representing its hybrid roots in LISP paradigms and XML syntax). 

While many contemporary engines utilized XML merely as a static data serialization format for layouts, the Useless engine utilized it as a fully homoiconic, functional programming language. KooLiXP natively supports advanced functional constructs such as \texttt{<map>}, \texttt{<fold>}, \texttt{<choose>}, and closures, allowing developers to define complex application logic entirely within the XML structure.

\subsection{The Compilation Pipeline: From Static AST to Executable Chunks}
By default, the engine's \texttt{XMLParser} reads the source files and generates a "static" (or dead) Abstract Syntax Tree (AST) in memory. Traversing this static tree sequentially (Interpreted Mode) is highly flexible but incurs a massive performance penalty due to continuous string lookups and tree parsing.

To bridge the gap between flexibility and high-performance execution, KooLiXP incorporates an Ahead-Of-Time (AOT) compiler. When the parser encounters a \texttt{<compile>} tag, the \texttt{XMLKompiler} is invoked. Using a highly decoupled static registration system (\texttt{XMLTagDict} and \texttt{XMLFactory}), the compiler maps each XML node to a specific C++ bytecode generator. 



The static XML tree is thereby transformed into an \textbf{Executable AST}---a linked graph of lightweight C++ \texttt{IChunk} objects. Because these chunks execute their logic via direct C++ virtual function calls (\texttt{IChunk::Execute(ExecutionState\&)}), the resulting execution speed is between $100\times$ and $1000\times$ faster than standard interpretation.

\subsection{Native C++ Method Bridging}
A critical requirement for KooLiXP was seamless integration with the underlying C++ engine. The \texttt{XMLChunksMethod.h} layer implemented a sophisticated template-based reflection system. Using \texttt{method\_adder} and \texttt{make\_method\_chunk}, native C++ class methods (e.g., \texttt{Widget::Resize}) could be bound directly to KooLiXP symbols. When a KooLiXP script invoked an action, the Executable AST executed a direct C++ member-function pointer without the overhead of runtime string evaluation.

\subsection{Compiler Optimizations}
During the transformation from static XML to executable \texttt{IChunk} graphs, the \texttt{XMLKompiler} applied several aggressive optimizations to maximize runtime throughput and minimize memory allocations:

\begin{itemize}
    \item \textbf{Execution Frame Compacting:} The KooLiXP Virtual Machine manages variable scope via an \texttt{ExecutionState} object. During compilation, the compiler performs static analysis on scope blocks (such as \texttt{<do>}). If a functional block does not declare any new local variables (\texttt{<let>}), the compiler entirely elides the creation of a new execution frame. Instead, it instructs the \texttt{IChunk} to safely reuse the parent's \texttt{ExecutionState}, drastically reducing stack depth and dynamic memory allocation during the engine's main loop.
    
    \item \textbf{Constant Folding and Evaluation:} In typical interpreted languages, declaring a literal value (e.g., \texttt{<integer value="5"/>}) forces the VM to allocate a temporary runtime variable. The \texttt{XMLKompiler} bypasses this by evaluating literals at compile-time. It generates specialized \texttt{CXML::EvaluableValue} chunks that embed the raw C++ primitives (\texttt{int}, \texttt{float}, \texttt{std::string}) directly into the executable AST, completely bypassing runtime variable assignments.
    
    \item \textbf{Symbol Interning and Direct Binding:} In typical interpreted environments, accessing a variable by name requires a relatively slow string-based hash-map lookup at runtime. The KooLiXP compiler optimizes this process using a two-tiered approach. First, string symbol names are interned using FourCC (Four-Character Code) representations. This converts the strings into 32-bit unsigned integers, reducing runtime symbol resolution to a highly efficient integer lookup. Second, depending on the specific chunk logic, the compiler often resolves symbols entirely during the compilation phase, storing a direct memory pointer to the corresponding child \texttt{IChunk} within the parent chunk. This completely eliminates the lookup overhead during the execution phase.
\end{itemize}

\section{Workspace Layering and Compositing}
Rather than treating the screen as a single flat canvas, the engine's \texttt{Workspace} class acts as a sophisticated compositor and event router. It explicitly partitions the screen into four distinct, Z-ordered \texttt{HubWidget} containers. This rigid stratification ensures that transient UI elements, like tooltips or context menus, never clip through or hide behind permanent application windows.



The \texttt{Workspace} manages the following distinct layers:
\begin{enumerate}
    \item \textbf{Desktop Layer (\texttt{\_desktop\_layer}):} The foundational \texttt{FormHub}. This layer contains persistent application frames, standard windows, and the primary application workspace. 
    \item \textbf{Popups Layer (\texttt{\_popups\_layer}):} Dedicated to dialogue boxes and secondary windows. This layer is managed cooperatively with the \texttt{PopupCycler} to handle both modal (blocking) and non-modal (floating) states.
    \item \textbf{Menus Layer (\texttt{\_menus\_layer}):} A high-priority transient layer for drop-down lists and context menus. It captures immediate focus upon activation.
    \item \textbf{Hints Layer (\texttt{\_hints\_layer}):} The absolute topmost visual layer, reserved for tooltips and hover-states. It is generally non-interactive but ensures critical user feedback is never obscured by active menus or overlapping windows.
\end{enumerate}

\subsection{Deterministic Flex-Style Layout System}
A standout feature of the engine's geometry layer is the \texttt{Layout} system. Long before CSS Flexbox became the standard for web browsers, the engine implemented a mathematically identical space-distribution algorithm. However, the design philosophy of the "Useless" layout engine makes it fundamentally superior to modern web compositing frameworks like React, Tailwind, or complex CSS.

Modern web compositing suffers from massive bloat. Developers must navigate hundreds of cascading rules, conflicting utility classes, DOM thrashing, and the heavy overhead of Virtual DOM reconciliation. The "Useless" layout engine avoids this completely by being fiercely single-purpose. It was designed to do layout well and keep it simple, entirely free from complex attributes or unnecessary web-bloat. 

\subsubsection{The Space Distribution Algorithm}
The layout engine does not parse massive style trees. Instead, the \texttt{Layout::Fill} algorithm operates directly on a simple array of widget boundaries, each defined by just four core mathematical properties: a base dimension ($d$), a flex weight ($w$), a minimum size ($min$), and a maximum size ($max$).

When a parent container is resized to a total dimension $D_{total}$, the algorithm first accounts for fixed padding (\texttt{padstart}, \texttt{padend}) and overlap. The remaining available space $S_{avail}$ is then distributed proportionally based on each widget's requested flex weight:

\begin{equation}
Size_{i} = d_{i} + S_{avail} \cdot \frac{w_{i}}{\sum w}
\end{equation}

\subsubsection{Multi-Pass Constraint Resolution}
What makes the system robust without being complicated is its built-in constraint solver. If allocating the proportional space causes a widget to exceed its $max$ bound (or fall below its $min$ bound), the algorithm instantly clamps the widget to its limit. The clamped widget's weight is then removed from the total pool, and the remaining space is seamlessly redistributed among the other flexible widgets in the same loop.

This pure C++ approach means there are no unpredictable "reflow" cascades or CSS specificity wars. A UI developer simply assigns a weight and boundaries, and the engine guarantees a perfect, deterministic layout in a mathematically proven resolution pass. By stripping away the bloated rulesets required by modern browsers, the engine achieves a layout pipeline that is faster, vastly more predictable, and significantly easier to reason about.

\subsection{Top-Down Event Routing and Modality}
Input events (mouse movements, clicks, and keyboard strokes) do not simply broadcast to all widgets simultaneously. The \texttt{Workspace} dispatches events strictly from the topmost layer downwards (Hints $\rightarrow$ Menus $\rightarrow$ Popups $\rightarrow$ Desktop). 

This top-down routing is the cornerstone of the engine's modal logic. The \texttt{Workspace} maintains an \texttt{\_input\_intercepting\_layer} state. If a widget in the Popups layer is flagged as modal, it invokes an \texttt{InputIntercept()} signal. The \texttt{Workspace} subsequently halts all downward event propagation at the Popups layer. The Desktop layer is effectively "starved" of input, instantly creating a robust modal lock without requiring the underlying widgets to query their active status.

\subsection{Graceful Focus Degradation}
To maintain a professional and predictable user experience, the \texttt{Workspace} and \texttt{PopupCycler} continuously monitor the \texttt{\_focus\_holding\_layer}. 

When a transient element completes its lifecycle---such as a user clicking outside of a context menu---the \texttt{Workspace::InvalidateFocus} method is triggered. The engine evaluates the remaining visible layers and performs a graceful focus degradation. If the Menus layer is emptied, focus automatically falls back to the Popups layer; if no popups are active, focus seamlessly returns to the Desktop layer. This automatic state cleanup prevents "dead focus" scenarios, ensuring the keyboard and mouse inputs are always routed to a valid target.

\section{Hardware Abstraction and System Integration}
The connection between the UI logic and the graphics hardware was orchestrated by the \texttt{GUIMaster}.

\subsection{System Bridge: GUIMaster}
The \texttt{GUIMaster} served as the central hub, initializing the \texttt{Screen} (hardware context), the \texttt{ScreenMan} (cursor and refresh manager), and the \texttt{Workspace}. It provided the \texttt{WidgetEnvironment} used by all components to access system resources like timers, clipboards, and soundcards.

\subsection{Multi-Backend Hardware Abstraction Layer (HAL)}
To ensure stable execution across the highly fragmented hardware landscape of the early 2000s, the engine completely decoupled its rendering logic from the physical graphics API. It accomplished this through unified \texttt{Screen} and \texttt{Surface} interfaces. 

The architecture officially supported four interchangeable backends:
\begin{enumerate}
    \item \textbf{OpenGL:} The primary cross-platform 3D hardware backend.
    \item \textbf{DirectDraw 7:} Optimized 2D hardware acceleration for legacy Windows machines.
    \item \textbf{GDI:} A pure software fallback for systems lacking any hardware acceleration.
    \item \textbf{RenderWare Graphics:} A proprietary backend utilizing Criterion's RenderWare (excluded from the open-source release).
\end{enumerate}

These backends could be used entirely interchangeably. Running a KooLiXP application on top of these drivers was conceptually identical to running a web application across different browsers. However, unlike modern web environments---which often suffer from bloat and inconsistent implementations---the Useless engine's abstraction guaranteed deterministic execution and 1:1 visual consistency across all platforms.



\subsubsection{The PixelTransfer Abstraction}
A major architectural hurdle in supporting multiple graphics APIs is the fundamental difference in how they manipulate video memory. 

APIs like OpenGL and GDI rely on functions that consume pixels from a provided system-memory buffer. Conversely, DirectDraw and RenderWare allow the application to explicitly lock the video surface and write bytes directly into VRAM. 

To hide this complexity from the application layer, the engine introduced the \texttt{PixelTransfer} abstract class. By utilizing \texttt{PixelTransfer} and its concrete implementations like \texttt{GenericTransfer}, the high-level UI logic never had to worry about the underlying routing of pixels. The engine automatically determined whether to push a buffer to an API or lock a surface and write directly, ensuring optimal memory throughput for the active backend.


\subsubsection{Surface Painters and the Specialized Pipeline}
Beyond simple blitting, the engine required a unified way to draw geometric primitives (lines, pixels, and polygons). Each backend provided a dedicated painter implementation, such as the \texttt{GLSurfacePainter} for OpenGL and the \texttt{GDISurfacePainter} for Windows GDI.

For APIs that permitted direct memory access (DirectDraw and RenderWare), the \texttt{Surface} implemented \texttt{LockableSurface}, and provided an optimized \texttt{SpecializedSurfacePainter}. Because locked surfaces could exist in numerous pixel formats (e.g., 16-bit 5-6-5, 24-bit RGB, or 32-bit ARGB), the \texttt{SpecializedSurfacePainter} used a \texttt{\_\_PIXEL\_SWITCH} macro. This macro dynamically selected the correct memory alignment and pixel format at runtime. It then routed the drawing commands to optimized pixel algorithms, such as a custom \texttt{DrawLine} routine, allowing the engine to render mathematically perfect, anti-aliased geometry directly into the locked buffer without relying on an external rasterization library.

\section{Conclusion}
The "Useless" engine was a milestone in vertically integrated architecture. By combining a functional compiled language (\textbf{KooLiXP}) with a layered compositing model and a multi-API hardware bridge, it provided a level of flexibility and performance that anticipated modern game engine standards. The modularity of its layering system, the intelligence of its ahead-of-time AST compiler, and the robustness of its focus management enabled the creation of highly polished, interactive interfaces that remained remarkably stable across vastly different hardware environments.

\end{document}